\documentclass[11pt]{article}
\usepackage[a4paper,margin=21mm]{geometry}
\usepackage{fontspec}
\setmainfont{Latin Modern Roman}
\newfontfamily\CJKfont{Noto Serif CJK SC}
\XeTeXlinebreaklocale "zh"
\XeTeXlinebreakskip=0pt plus 1pt
\usepackage{amsmath,amssymb,amsthm,amscd,graphicx,color}
\usepackage{xurl}
\usepackage[unicode,hidelinks]{hyperref}
\allowdisplaybreaks
\setlength\emergencystretch{3em}
\numberwithin{equation}{section}

\usepackage[all]{xy}

\numberwithin{equation}{section}

\newtheorem{Theorem}{Theorem}[section]
\newtheorem*{Theorem*}{Theorem}
\newtheorem{Corollary}[Theorem]{Corollary}
\newtheorem{Lemma}[Theorem]{Lemma}
\newtheorem{Claim}[Theorem]{Claim}
\newtheorem{Conjecture}[Theorem]{Conjecture}
\newtheorem{Proposition}[Theorem]{Proposition}
 { \theoremstyle{definition}
\newtheorem{Definition}[Theorem]{Definition}
\newtheorem{Note}[Theorem]{Note}
\newtheorem{Example}[Theorem]{Example}
\newtheorem{Remark}[Theorem]{Remark} }

\newcommand{\rk}{\operatorname{rk}}	
\newcommand{\Quot}{\mathsf{Quot}}
\newcommand{\sym}{\textnormal{Sym}}

\begin{document}
\begin{center}\Large Big and nef tautological bundles from very ample line bundles on minimal general-type surfaces under the stated numerical and even-exponent conditions\end{center}
\noindent\textbf{Stable ID:} 1251\quad\textbf{Author:} Dragos Oprea\par
\noindent\textbf{Work:} Big and Nef Tautological Vector Bundles over the Hilbert Scheme of Points\par
\noindent\textbf{Source:} \url{https://sigma-journal.com/2022/061/}\par
\noindent\textbf{Version:} Official SIGMA 18 (2022) article 061, DOI 10.3842/SIGMA.2022.061; journal PDF and native TeX acquired 2026-09-30, exact bytes pinned by SHA256\par
\noindent\textbf{Rights:} CC BY 4.0. Authors retain copyright. \url{https://creativecommons.org/licenses/by/4.0/}. Reuse and typesetting adaptation; original language and mathematical content preserved.\par
\noindent\textbf{Difficulty (editorial):} {\CJKfont H2: The proof rewrites the top Segre integral using the established universal generating series and then controls its sign through two nontrivial residue substitutions. In the fully expanded even-exponent case, these substitutions separate a positive range from a vanishing coefficient range, allowing a strict positivity argument after convolution. The surface inequalities then guarantee a positive top self-intersection of the globally generated tautological bundle; simply citing the universal Segre formula would not supply that sign.}\par
\medskip\noindent\textbf{Editorial boundary note (not author text):} Selected Theorem4.4 for a smooth projective minimal general-type surface, positive k, (k-1)-very ample line bundle L with chi(L)>=3k and L.K>=2K\textasciicircum{}2+k+1, in the previously reviewed explicit all-even/nonnegative-n branch: m=L.K-2K\textasciicircum{}2, n=(L-K)\textasciicircum{}2+3chi(OX)-4k-1, p=K\textasciicircum{}2-chi(OX)+3 all even and n>=0. Full printed broader theorem and Lemma4.1 remain literal context; the other parity branches whose details the author leaves to the reader are unselected. Complete Claim4.2 double residue-change/coefficient convolution, full even-branch positivity and main numerical/Noether reduction are retained. The exact earlier4.2 residue substitution is included with its original K3-blowup setting as context; that K3 outcome is unselected. Universal Segre formula and established globally-generated/top-Segre big-nef criterion stay author-cited prerequisites. Every source index/sign/exponent is unchanged, including printed summation-variable wording; no new proof or generalized parity completion.\par
\medskip\hrule\medskip\noindent\textbf{Author text begins below.}\par
\setcounter{section}{0}
% Original source lines 63--67
\section{Introduction}Let $X$ be a smooth projective surface, and let $X^{[k]}$ denote the Hilbert scheme of $k$ points on $X$. Each vector bundle $F\to X$ of rank $r$ yields a tautological vector bundle $F^{[k]}\to X^{[k]}$ of rank $rk$ given by \vspace{-.5ex}
\begin{gather*}
F^{[k]}=p_{\star} (q^{\star} F\otimes \mathcal O_{\mathcal Z}).
\end{gather*}
Here, $p$, $q$ are the natural projections from $X^{[k]}\times X$, and $\mathcal Z\subset X^{[k]}\times X$ denotes the universal subscheme.


% Original source lines 79--84
Recall that a vector bundle $V$ over a scheme $Y$ is said to be big and nef if the line bundle $\mathcal O_{\mathbb P(V)}(1)\to \mathbb P(V)$ is big and nef, where $\mathbb P(V)$ denotes the projective bundle of one dimensional quotients. A discussion of big and nef vector bundles can be found in \cite[Chapters~6 and~7]{L-II}.
A useful well-known characterization occurs when $V\to Y$ is globally generated. In this case, if the top Segre class
\begin{equation}
\label{pos}(-1)^{\dim Y} \int_{Y} s(V)>0
\end{equation}
it follows that $V\to Y$ is big and nef.

\setcounter{section}{1}
% Original source lines 122--130
\section[K3 surfaces]{$\boldsymbol{K3}$ surfaces} \label{seck3} Let $X$ be a smooth projective surface. The bundle $F^{[k]}\to X^{[k]}$ is globally generated, and therefore nef, provided $F\to X$ is $(k-1)$-very ample. By definition, $(k-1)$-very ampleness is the requirement that the natural map
\begin{gather*}
H^0(X, F)\to H^0(X, F\otimes \mathcal O_{\zeta})
\end{gather*}
is surjective for all zero-dimensional subschemes $\zeta$ of $X$ of length $k$. Thus, via \eqref{pos}, if $F$ is $(k-1)$-very ample and
\begin{gather*}
\int_{X^{[k]}} s\big(F^{[k]}\big)>0,
\end{gather*}
then $F^{[k]}\to X^{[k]}$ is big and nef. This is explained for instance in \cite[Propositions~2.4 and~4.5]{BBF}.

\setcounter{section}{3}
% Original source lines 536--696
\section{Other geometries} \label{secother}
The Segre integrals for arbitrary surfaces are established only when $F$ has rank $1$ or $2$, see \cite{MOP1, MOP}. In rank $1$, the answers were conjectured by Lehn \cite{Le}. The formulation below can be found in \cite{MOP}:
\begin{equation}\label{lehnc}
\sum_{n=0}^{\infty} z^k \int_{X^{[k]}} s\big(L^{[k]}\big) = A_1 (z)^{L^2} A_2(z)^{\chi(\mathcal O_X)} A_3 (z)^{L.K_X} A_4(z)^{K_X^2}.
\end{equation}
Here, for $z=t(1+2t)^2$, we have
\begin{gather*}
A_1(z)=(1+2t)^{\frac{1}{2}},
\\
A_2(z)=(1+2t)^{\frac{3}{2}} (1+6t)^{-\frac{1}{2}},
\\
A_3(z)=\frac{1}{2} (1+2t)^{-1} \left(\sqrt{1+2t}+\sqrt{1+6t}\right),
\\
 A_4(z)=4 (1+2t)^{\frac{1}{2}} (1+6t)^{\frac{1}{2}} \left(\sqrt{1+2t}+\sqrt{1+6t}\right)^{-2}.
 \end{gather*}

\subsection{A positivity result} It is more difficult to determine the sign of the top Segre integral from the above formulas. The~following lemma plays an important role in the analysis. We will apply it in the next section to geometric situations.

\begin{Lemma}\label{cl}
For all integers $m$, $n$, $p$ with $m\geq 0$, $p\geq 0$ such that $m+n+p$ is even, the series
\begin{gather*}
f(t)=\left(\sqrt{1+2t}+\sqrt{1+6t}\right)^m {(1+2t)}^{\frac{n-1}{2}} (1+6t)^{\frac{p-1}{2}}
\end{gather*}
has positive coefficients up to order less or equal than $\min\big(\frac{1}{2}(m+n+p)-1, m-1\big)$.
\end{Lemma}
Taking the minimum is necessary. Indeed, for $(m, n, p)=(2, 19, 1)$ the term of order $\frac{1}{2}(m+n+p)-1$ has negative coefficient, while for $(m, n, p)=(4, 0, 0)$, the term of order $m-1$ has zero coefficient. The lemma also holds for $m+n+p$ odd, but this case {\it never} occurs geometrically. The hypothesis $m, p\geq 0$ can fail in geometric examples, but our proof requires it.

\begin{proof}
The argument is not straightforward due to the alternating signs
in the expansions
\begin{gather*}
\sqrt{1+2t}= 1+t - \frac{t^2}{2}+\frac{t^3}{2}-\frac{5t^4}{8}+\cdots,
\\
\sqrt{1+6t}=1 + 3t - \frac{9t^2}{2}+\frac{27t^3}{2} - \frac{405 t^4}{8} +\cdots.
\end{gather*}
However, a residue calculation and suitable changes of variables will render the answer manifestly positive.

We begin by considering the case $n\geq 0$. In this situation, assume first $m$, $n$, $p$ are even, and~set
\begin{gather*}
F(t)=\left(\sqrt{1+2t}+\sqrt{1+6t}\right)^m {(1+2t)}^{-\frac{1}{2}} (1+6t)^{-\frac{1}{2}}.
\end{gather*}

\begin{Claim}\label{claim}
The series $F$ has positive coefficients up to order less or equal than $\frac{m}{2}-1$, and no terms of order between $\frac{m}{2}$ and $m-1$.
\end{Claim}


We assume this for now. Note that
\begin{gather*}
f(t)=F(t) P(t), \quad P(t)=(1+2t)^{\frac{n}{2}} (1+6t)^{\frac{p}{2}}.
\end{gather*}
Clearly, $P$ is a polynomial of degree $\frac{n+p}{2}$ with positive coefficients. This observation together with Claim~\ref{claim} gives the argument. Indeed, let $k\leq \min\big(\frac{1}{2}(m+n+p)-1, m-1\big)$. Writing $F_i$ and $P_j$ for the coefficients of $F$ and $P$, we have
\begin{gather*}
\operatorname{Coeff}_{ t^k} f(t)=\sum F_iP_j, \qquad\text{the sum ranging over}\quad i+j=k, \quad i\geq 0,\quad 0\leq j\leq \frac{n+p}{2}.
\end{gather*}
The sum is nonnegative since $i\leq k\leq m-1$, so $F_i\geq 0$ by the Claim~\ref{claim}, and $P_j> 0$. The sum is in fact strictly positive. Indeed, for $k\geq \frac{n+p}{2}$, it contains the term
\begin{gather*}
F_{k-\frac{n+p}{2}} P_{\frac{n+p}{2}}>0.
\end{gather*}
The last statement is also a consequence of Claim~\ref{claim}, using that $k-\frac{n+p}{2}\leq\frac{m}{2}-1$. When $k<\frac{n+p}{2}$, the sum contains the term $F_0 P_k=2^m P_k>0$.

\begin{proof}[Proof of Claim~\ref{claim}]
We seek to show that for all $k\leq \frac{m}{2}-1$, the residue
\begin{gather*}
\operatorname{Res}_{t=0} \frac{F(t)}{t^{k+1}} \,{\mathrm dt}>0.
\end{gather*}
The peculiar change of variables
\begin{gather*}
t=\frac{s(s+1)}{2\sqrt{3}s+(2+\sqrt{3})}
\end{gather*}
will simplify the calculation.
For convenience, set
\begin{gather*}
a=\frac{2-\sqrt{3}}{2\sqrt{3}}>0\qquad\text{so that}\quad t=\frac{s(s+1)}{2\sqrt{3}(s+a+1)}.
\end{gather*}
We note the following identities
\begin{gather*}
\sqrt{1+2t}=\frac{s+\frac{\sqrt{3}+1}{2}}{3^{\frac{1}{4}} (s+a+1)^{\frac{1}{2}}}, \qquad \sqrt{1+6t}=\frac{3^{\frac{1}{4}}\big(s+\frac{\sqrt{3}+1}{2\sqrt{3}}\big)}{(s+a+1)^{\frac{1}{2}}},
\\
\sqrt{1+2t}+\sqrt{1+6t}=\frac{(\sqrt{3}+1)(s+1)}{3^{\frac{1}{4}}(s+a+1)^{\frac{1}{2}}},\qquad
{\mathrm dt}=\frac{\big(s+\frac{\sqrt{3}+1}{2}\big)\big(s+\frac{\sqrt{3}+1}{2\sqrt{3}}\big)}{2\sqrt{3} (s+a+1)^2}\,{\mathrm ds}.
\end{gather*}
From here, we obtain
\begin{gather*}
\operatorname{Res}_{t=0} \frac{F(t)}{t^{k+1}} \,{\mathrm dt}=\operatorname{Res}_{s=0} \frac{G(s)}{s^{k+1}} \,{\mathrm ds},
\end{gather*}
where
\begin{gather*}
G(s)=\alpha (s+1)^{m-k-1} (s+a+1)^{k-\frac{m}{2}}, \qquad
\alpha=2^k \big(\sqrt{3}+1\big)^{m} \sqrt{3}^{k-\frac{m}{2}}>0.
\end{gather*}
A second change of variables will be needed next (this change of variables could have been carried out simultaneously with the first, but this is a price worth paying for readability). We~set
\begin{gather*}
s=\frac{(a+1)^2}{a} \frac{u}{1-\frac{a+1}{a}u}.
\end{gather*}
The reader can verify that
\begin{gather*}
s+1=\frac{1+(a+1)u}{1-\frac{a+1}{a}u}, \qquad s+a+1=\frac{a+1}{1-\frac{a+1}{a}u}, \qquad
{\mathrm ds}=\frac{(a+1)^2}{a} \frac{1}{\big(1-\frac{a+1}{a}u\big)^2}\,{\mathrm du}.
\end{gather*}
By direct calculation, we find
\begin{gather*}
\operatorname{Res}_{s=0} \frac{G(s)}{s^{k+1}}\,{\mathrm ds}=\operatorname{Res}_{u=0} \frac{H(u)}{u^{k+1}}\,{\mathrm du}
\end{gather*}
for
\begin{gather*}
H(u)=\beta (1+(a+1)u)^{m-k-1} \biggl(1-\frac{a+1}{a}u\biggr)^{k-\frac{m}{2}}, \qquad \beta=\alpha a^{k} (a+1)^{-k-\frac{m}{2}}>0.
\end{gather*}
The first term is a polynomial with positive coefficients $a_i$ given by binomial numbers, for $0\leq i\leq m-k-1$. The second term also has positive coefficients
\begin{gather*}
b_j=\biggl(-\frac{a+1}{a}\biggr)^j\binom{k-\frac{m}{2}}{j}>0
\end{gather*}
since $k-\frac{m}{2}<0$. Thus
\begin{gather*}
\operatorname{Coeff}_{u^k} H(u)=\!\sum a_i b_j,\quad \text{the sum ranging over}\!\quad i+j=k, \quad 0\leq i\leq m-k-1, \quad j\geq 0.
\end{gather*}
This coefficient is positive since $a_i>0$, $b_j> 0$ and the sum is non-empty (the term $a_0b_k=b_k>0$ appears in the sum).

When $m$ is even, and $\frac{m}{2}\leq k\leq m-1$, the expression $H(u)$ is a polynomial in $u$ of degree $\frac{m}{2}-1<k$. Hence, the coefficient of $u^k$ in $H(u)$ vanishes.
This proves the claim (and completes the argument when $m$, $n$, $p$ are even and $n\geq 0$).
\end{proof}

When $n\geq 0$, there are three other cases to consider, which require different choices for $F$ and $P$:
\begin{itemize}\itemsep=0pt
\item [$(i)$] when $m$ even, $n$, $p$ odd, we set
\begin{gather*}
F(t)=\left(\sqrt{1+2t}+\sqrt{1+6t}\right)^m, \qquad
P(t)=(1+2t)^{\frac{n-1}{2}} (1+6t)^{\frac{p-1}{2}},
\end{gather*}
\item [$(ii)$] when $m$ odd, $n$ even, $p$ odd, we set
\begin{gather*}
F(t)=\left(\sqrt{1+2t}+\sqrt{1+6t}\right)^m {(1+2t)}^{-\frac{1}{2}}, \qquad
P(t)=(1+2t)^{\frac{n}{2}} (1+6t)^{\frac{p-1}{2}},
\end{gather*}
\item [$(iii)$] when $m$ odd, $n$ odd, $p$ even, we set
\begin{gather*}
F(t)=\left(\sqrt{1+2t}+\sqrt{1+6t}\right)^m (1+6t)^{-\frac{1}{2}}, \qquad
P(t)=(1+2t)^{\frac{n-1}{2}} (1+6t)^{\frac{p}{2}}.
\end{gather*}
\end{itemize}
In case $(i)$, $F$ has positive coefficients up to order $\frac{m}{2}$, and no terms of order between $\frac{m}{2}+1$ and $m-1$.
In cases $(ii)$ and~$(iii)$, $F$ has positive coefficients up to order $\frac{m-1}{2}$, and no terms of order between $\frac{m+1}{2}$ and $m-1$. The arguments are similar, and we leave the details to the reader.

When $n<0$, the reasoning used to prove Claim~\ref{claim} also works to deal directly with the function
\begin{gather*}
f(t)=\left(\sqrt{1+2t}+\sqrt{1+6t}\right)^m {(1+2t)}^{\frac{n-1}{2}} (1+6t)^{\frac{p-1}{2}}.
\end{gather*}
Following exactly the same steps, first changing from $t$ to $s$ and then from $s$ to $u$, we obtain
\begin{align*}
H(u)={}&\gamma (1+(a+1)u)^{m-k-1} \biggl(1-\frac{a+1}{a}u\biggr)^{k-\frac{m+n+p}{2}} \biggl(1-\frac{(\sqrt{3}+1)^3}{4\sqrt{3}}u\biggr)^{n}
\\
&\times \biggl(1+\frac{\big(\sqrt{3}+1\big)^3}{4}u\biggr)^p,
\end{align*}
for $\gamma>0$. By the same arguments, this has positive coefficients when the exponents
\begin{gather*}
m-k-1\geq 0, \qquad k-\frac{m+n+p}{2}<0, \qquad n<0, \qquad p\geq 0,
\end{gather*}
which we assumed.
\end{proof}




% Original source lines 697--714
\subsection[K3 blowups]{$\boldsymbol{K3}$ blowups}
The top Segre classes computed by formula \eqref{lehnc} are always coefficients in series of the form~$f(t)$ as in Lemma \ref{cl}, barring the condition $m, p\geq 0$. When this condition is satisfied, we easily obtain big and nef criteria for the tautological bundles $L^{[k]}\to X^{[k]}$.

There are several specific examples where our techniques apply. We illustrate them first when
\begin{gather*}
\pi\colon\ X\to S
\end{gather*}
is the blowup of a $K3$ surface $S$ at a point $p\in S$. We assume
$S$ has Picard rank $1$, with ample Picard generator $H$. Let $H^2=2h$.

\begin{Theorem}\label{blok}
Let $E$ be the exceptional divisor on $X$, and set $L=H-\ell E$. Assume $\ell\geq k-1$ and
\begin{gather*}
2h>\max \bigl((\ell+2)^2-6, (\ell+1)^2+4k, \ell(\ell+1)+6k-6\bigr).
\end{gather*}
Then $L^{[k]}$ is big and nef on $X^{[k]}$.
\end{Theorem}



% Original source lines 778--791
It remains to explain that the top Segre class of $L^{[k]}$ is positive. We use \eqref{lehnc} and we change variables
\begin{gather*}
z=t(1+2t)^2\qquad \text{so that}\quad \mathrm dz=(1+2t)(1+6t)\,{\mathrm dt}.
\end{gather*}
Then
\begin{align}
\int_{X^{[k]}} s\big(L^{[k]}\big)&=\operatorname{Res}_{z=0} A_1^{2h-\ell^2} A_2^2 A_3^{\ell} A_4^{-1} \frac{\mathrm dz}{z^{k+1}}\nonumber
\\
\nonumber&=
\operatorname{Res}_{t=0} 2^{-\ell-2} (1+2t)^{h-\frac{\ell^2}{2}-2k-\ell+\frac{3}{2}} (1+6t)^{-\frac{1}{2}} \left(\sqrt{1+2t}+\sqrt{1+6t}\right)^{\ell+2}\frac{\mathrm dt}{t^{k+1}}
\\
&= 2^{-\ell-2} \operatorname{Coeff}_{t^k} (1\!+2t)^{h-\frac{\ell^2}{2}-2k-\ell+\frac{3}{2}} (1\!+6t)^{-\frac{1}{2}} \left(\sqrt{1+2t}\!+\sqrt{1+6t}\right)^{\ell+2}\!.
\label{cvvv}
\end{align}


% Original source lines 799--832
\subsection{Surfaces of general type} It is natural to wonder how far these techniques take us. We show

\begin{Theorem} \label{gtyp}Assume $X$ is a smooth projective minimal surface of general type. Let $L$ be a~$(k-1)$-very ample line bundle such that
\begin{gather*}
\chi(L)\geq 3k, \qquad L.K_X\geq 2K_X^2+k+1.
\end{gather*}
Then $L^{[k]}$ is big and nef over $X^{[k]}$.
\end{Theorem}

\begin{proof}
Arguing as in the proof of Theorem \ref{blok}, in particular equation \eqref{cvvv}, we express the Segre integral as the $t^k$-coefficient in the series
\begin{gather*}
f(t)=2^{-m} \left(\sqrt{1+2t}+\sqrt{1+6t}\right)^m {(1+2t)}^{\frac{n-1}{2}} (1+6t)^{\frac{p-1}{2}},
\end{gather*}
where
\begin{gather*}
m=L.K-2K^2, \qquad n=(L-K)^2+3\chi - 4k - 1, \qquad p=K^2-\chi+3,
\end{gather*}
with $K=K_X$ and $\chi=\chi(\mathcal O_X)$. Note that
\begin{gather*}
m+n+p=2\chi(L)-4k+2
\end{gather*}
is even. Furthermore, $p\geq 0$. Indeed, let $p_g$, $q$ denote the genus and irregularity of $X$. By~Noet\-her's inequality $K^2\geq 2p_g-4$, and $K^2>0$ since $X$ is minimal of general type. Averaging, we obtain $K^2>p_g-2$, and thus
\begin{gather*}
p=K^2-\chi+3>(p_g-2)-(1-q+p_g)+3=q\geq 0.
\end{gather*}
By Lemma \ref{cl}, the coefficients of
$f(t)$ up to order $\min\big(\frac{1}{2} (m+n+p)-1, m-1\big)$ are positive. In~particular, if
\begin{gather*}
k\leq \min\bigg(\frac{1}{2} (m+n+p)-1, m-1\bigg),
\end{gather*}
then the Segre integral is positive. The latter inequality is exactly our hypothesis, and thus $L^{[k]}$ is big and nef.
\end{proof}

\begin{thebibliography}{99}
\bibitem[7]{BBF}
Bini G., Boissi\`ere S., Flamini F., Some families of big and stable bundles on
 $K3$ surfaces and on their Hilbert schemes of points, \href{https://arxiv.org/abs/2109.01598}{arXiv:2109.01598}.

\bibitem[21]{L-II}
Lazarsfeld R., Positivity in algebraic geometry. {II}.Positivity for vector
 bundles, and multiplier ideals, \textit{Ergebnisse der Mathematik und ihrer
 Grenzgebiete. 3.~Folge. A Series of Modern Surveys in Mathematics}, Vol.~49,
 \href{https://doi.org/10.1007/978-3-642-18810-7}{Springer-Verlag}, Berlin, 2004.

\bibitem[22]{Le}
Lehn M., Chern classes of tautological sheaves on {H}ilbert schemes of points
 on surfaces, \href{https://doi.org/10.1007/s002220050307}{\textit{Invent. Math.}} \textbf{136} (1999), 157--207,
 \href{https://arxiv.org/abs/math.AG/9803091}{arXiv:math.AG/9803091}.

\bibitem[26]{MOP1}
Marian A., Oprea D., Pandharipande R., The combinatorics of {L}ehn's
 conjecture, \href{https://doi.org/10.2969/jmsj/78747874}{\textit{J.~Math. Soc. Japan}} \textbf{71} (2019), 299--308,
 \href{https://arxiv.org/abs/1708.08129}{arXiv:1708.08129}.

\bibitem[27]{MOP}
Marian A., Oprea D., Pandharipande R., Higher rank {S}egre integrals over the
 {H}ilbert scheme of points, \href{https://doi.org/10.4171/jems/1149}{\textit{J.~Eur. Math. Soc. (JEMS)}} \textbf{24}
 (2022), 2979--3015, \href{https://arxiv.org/abs/1712.02382}{arXiv:1712.02382}.

\end{thebibliography}
\end{document}
